% from ega0-1-fr.tex
\oldpage[0\textsubscript{I}]{51}

ces homomorphismes (pour $n$ variable) définissent un
\emph{di-homomorphisme de modules gradués}
\[
  \Gamma_*(\mathcal{L},\mathcal{F})
  \longrightarrow
  \Gamma_*(f^*(\mathcal{L}),f^*(\mathcal{F})).
\]
\end{env}

\begin{env}[5.4.7]
\EGAReaderTarget{canon-ref.0.5.4.7-fr}{0.5.4.7-fr}
On peut montrer qu'il existe un \emph{ensemble} $\mathfrak{M}$ (noté aussi
$\mathfrak{M}(X)$) de $\mathcal{O}_X$-Modules inversibles tel que tout
$\mathcal{O}_X$-Module inversible soit isomorphe à un élément de
$\mathfrak{M}$ et un seul\footnote{Voir le livre en préparation cité dans
l'Introduction.}~; on définit dans $\mathfrak{M}$ une loi de composition en
faisant correspondre à deux éléments $\mathcal{L}$, $\mathcal{L}'$ de
$\mathfrak{M}$ l'unique élément de $\mathfrak{M}$ isomorphe à
$\mathcal{L}\otimes_{\mathcal{O}_X}\mathcal{L}'$. Avec cette loi de
composition, $\mathfrak{M}$ est un \emph{groupe isomorphe au groupe de
cohomologie} $H^1(X,\mathcal{O}_X^*)$, où $\mathcal{O}_X^*$ est le
sous-faisceau de $\mathcal{O}_X$ tel que $\Gamma(U,\mathcal{O}_X^*)$ soit le
groupe des éléments inversibles de l'anneau $\Gamma(U,\mathcal{O}_X)$ pour
tout ouvert $U\subset X$ ($\mathcal{O}_X^*$ est donc un faisceau de groupes
abéliens \emph{multiplicatifs}).

On notera pour cela que pour tout ouvert $U\subset X$, le groupe des sections
$\Gamma(U,\mathcal{O}_X^*)$ s'identifie canoniquement au \emph{groupe des
automorphismes} du $(\mathcal{O}_X|U)$-Module $\mathcal{O}_X|U$,
l'identification faisant correspondre à une section $\varepsilon$ de
$\mathcal{O}_X^*$ au-dessus de $U$ l'automorphisme $u$ de $\mathcal{O}_X|U$
tel que $u_x(s_x)=\varepsilon_xs_x$ pour tout $x\in X$ et tout
$s_x\in\mathcal{O}_x$. Soit alors $\mathfrak{U}=(U_\lambda)$ un recouvrement
ouvert de $X$~; la donnée, pour tout couple d'indices $(\lambda,\mu)$ d'un
automorphisme $\theta_{\lambda\mu}$ de
$\mathcal{O}_X|(U_\lambda\cap U_\mu)$ revient à se donner une
$1$-\emph{cochaîne} du recouvrement $\mathfrak{U}$, à valeurs dans
$\mathcal{O}_X^*$, et dire que les $\theta_{\lambda\mu}$ vérifient la
condition de recollement (3.3.1) signifie que la cochaîne correspondante est
un \emph{cocycle}. De même, la donnée, pour tout $\lambda$, d'un automorphisme
$\omega_\lambda$ de $\mathcal{O}_X|U_\lambda$ revient à la donnée d'une
$0$-cochaîne du recouvrement $\mathfrak{U}$, à valeurs dans
$\mathcal{O}_X^*$, et son \emph{cobord} correspond à la famille des
automorphismes
$(\omega_\lambda|U_\lambda\cap U_\mu)\circ
 (\omega_\mu|U_\lambda\cap U_\mu)^{-1}$.

On peut faire correspondre à tout $1$-cocycle de $\mathfrak{U}$ à valeurs
dans $\mathcal{O}_X^*$ l'élément de $\mathfrak{M}$ isomorphe au
$\mathcal{O}_X$-Module inversible obtenu par recollement à partir de la
famille d'automorphismes $(\theta_{\lambda\mu})$ correspondant à ce cocycle,
et à deux cocycles cohomologues correspondent deux éléments égaux de
$\mathfrak{M}$ (3.3.2)~; autrement dit, on a ainsi défini une application
$\varphi_{\mathfrak{U}}:H^1(\mathfrak{U},\mathcal{O}_X^*)\to\mathfrak{M}$.
En outre, si $\mathfrak{B}$ est un second recouvrement ouvert de $X$, plus fin
que $\mathfrak{U}$, le diagramme
\[
\begin{tikzcd}[column sep=huge,row sep=large]
H^1(\mathfrak{U},\mathcal{O}_X^*)
  \arrow[dd]
  \arrow[rrd,"\varphi_{\mathfrak{U}}"] & & \\
& & \mathfrak{M} \\
H^1(\mathfrak{B},\mathcal{O}_X^*)
  \arrow[rru,swap,"\varphi_{\mathfrak{B}}"] & &
\end{tikzcd}
\]
où la flèche verticale est l'homomorphisme canonique (G, II, 5.7), est
commutatif, comme il résulte de (3.3.3). Par passage à la limite inductive, on
obtient donc bien une application
$H^1(X,\mathcal{O}_X^*)\to\mathfrak{M}$, le groupe de cohomologie de \v{C}ech
$\check{H}^1(X,\mathcal{O}_X^*)$ s'identifiant comme on sait au premier groupe
de cohomologie $H^1(X,\mathcal{O}_X^*)$ (G, II, 5.9, cor. du th.~5.9.1).
Cette application est \emph{surjective}~: en effet, par définition, pour tout
$\mathcal{O}_X$-Module inversible $\mathcal{L}$, il y a un recouvrement ouvert
$(U_\lambda)$ de $X$ tel que $\mathcal{L}$ s'obtienne par recollement des
faisceaux $\mathcal{O}_X|U_\lambda$ (3.3.1). Elle est aussi
\emph{injective}, car il suffit de le prouver pour les applications
$H^1(\mathfrak{U},\mathcal{O}_X)\to\mathfrak{M}$, et cela résulte alors de
(3.3.2). Il reste à montrer que

